\begin{examplebox}{Example 3 --- Ecological Pyramids}
\textbf{Student writing} (Score 3):

\begin{quote}\itshape
Ecologists construct a pyramid of numbers based on the population sizes of organisms in each trophic level. Sometimes population sizes decrease at each higher trophic level. This is not always true. For example, one tree can be food for 50,000 insects. In this case, the pyramid would be inverted.

Energy decreases at each succeeding trophic level. The total energy transfer from one trophic level to the next is only about ten percent because organisms fail to capture and eat all the food available at the trophic level below them. When an organism consumes food, it uses some of the energy in the food for metabolism, some for building body tissues, and some is given off as waste.
\end{quote}

\textbf{GM rewrite} (sim $0.75$):

\begin{quote}\itshape
Ecologists construct a pyramid of numbers based on the population sizes of organisms in each trophic level. Sometimes population sizes decrease at each higher trophic level. This is not always true. For example, one tree can be food for 50,000 insects. In this case, the pyramid would be inverted. Energy decreases at each succeeding trophic level. The total energy transfer from one trophic level to the next is only about ten percent because organisms fail to capture and eat all the food available at the trophic level below them. When an organism consumes food, it uses some of the energy in the food for metabolism, some for building body tissues, and some is given off as waste. When the organism is eaten, the energy that was used to build body tissue is available as energy to be used by the organism that consumed it. The energy lost at each successive trophic level enters the environment as heat.
\end{quote}

\textbf{ME rewrite} (sim $0.34$):

\begin{quote}\itshape
Ecologists use ecological pyramids to model the distribution of matter and energy within an ecosystem. This type of pyramid shows how energy flows through an ecosystem. The initial source of energy for all ecological pyramids is energy from the sun, which decreases at each succeeding trophic level due to inefficient energy transfer. The total energy transfer from one trophic level to the next is only about ten percent. When an organism consumes food, it uses some of the energy in the food for metabolism, some for building body tissues, and some is given off as waste, with the energy lost at each successive trophic level entering the environment as heat. Ecologists also construct pyramids of numbers based on population sizes, which sometimes decrease at each higher trophic level, but can also be inverted, such as when one tree is food for 50,000 insects.
\end{quote}
\end{examplebox}
